% TJ93090C
% Pneumothorax
% 14.0
% 2013-11-30
\label{m:pneumothorax}
\EE{Taktik}: \Speech{\Ca{Zeitkritisch! Vitale Bedrohung!}
Entlastungspunktion, zügiger Transport
in eine geeignete Einrichtung}

\begin{itemize}
  \item \TxMassVitMKBes
  \begin{itemize}
    \item \EE{Lagerung}: Oberkörper hoch. Wenn stabile
    Seitenlage notwendig ist, dann Lagerung auf die
    \E{verletzte} Seite, damit die unverletzte
    Lunge  nicht unnötig belastet wird.
    \item \Z{\EE{Notarzt}: Ein Pneumothorax kann mittels
    Drainage bzw. Punktion entlastet werden (ärztliche Maßnahme).}
  \end{itemize}
  \item \EE{Verband} bei offener Thoraxverletzung: \EE{Nicht
  luftdicht verbinden}, es kann sonst ein
  Spannungspneumothorax entstehen! Steril abdecken.
  \item \EE{Transportentscheidung}: Schockraum
\end{itemize}